\documentclass[10pt]{article}

\usepackage[utf8]{inputenc}

\usepackage[T1]{fontenc}

\usepackage{amsmath}

\usepackage{amsfonts}

\usepackage{amssymb}

\usepackage[version=4]{mhchem}

\usepackage{stmaryrd}

\usepackage{bbold}

\usepackage{mathrsfs}

\usepackage{graphicx}

\usepackage[export]{adjustbox}

\graphicspath{ {./images/} }



\begin{document}

Лит.: [1] К р а с н о с е л ь с к и й М. А., Топологические методы в теории нелинейных интегральных уравнений, М , 1956 , [2] Интег ральные уравнения, М., 1968 (Справочная матем б-ка)



УРЫСОНА - БРАУРРА ЛЕММА Б. В. Хведелидзе. Б р а у э а - Т ц е ле м м а,- утверждение о возможности продолжения непрерывных функций $c$ Іодпространства топологич. пространства на все пространство. Пусть $X$ - нормальное пространство и $F$ єго замкнутое подмножество. Тогда любую непрерывную функцию $f: F \rightarrow \mathbb{R}$ можно продолжить непрерывно до функции $g: X \rightarrow \mathbb{R}$, т. е. можно найти такую нешрерывную функцию $g$, что $g(x)=f(x)$ для всех $x \in F$. При этом если функция $f$ ограничена, то существует такое ее продолжение $g$, что



\[

\sup _{x \in F}|f(x)|=\sup _{x \in X}|g(x)|

\]



У.- Б. л. была доказана Л. Брауэром (L. Brouwer) и А. Лебегом (Н. Lebesgue) для $X=\mathbb{R}^{n}$, А. Тице (A. Titze) для произвольного метрич. пространства $X$ и П. С. Урысоном - в приведенной выше формулировке (к-рая может служить характеризацией нормальных пространств и является, таким образом, окончателььної).



Jum.: [1] у р ы с о н П. С., "Math. Ann.», 1925, Bd 94, и другим областям математики, т. 1 , - J 1951 , c. 177-218)



М. Г. Колевникова.



УСЕЧЕННОЕ РАСПРЕДЕЛЕНИЕ - распределение вероятностеӥ, получаемое из данного распределения перенесением массы, заключенной вне нек-рого фиксированного отрезка, на этот отрезок. Пусть вероятностное распределение на прямой задано функцией распределения $F(x)$. У с е че н ны м р а с п р е д е л ен и е м, отвечающим $F$, наз. распределение с функцией распределения



\[

F_{a, b}(x)=\left\{\begin{array}{cll}

0 & \text { при } & x \leqslant a, \\

\frac{F(x)-F(a)}{F(b)-F(a)} & \text { при } & a<x \leqslant b, \\

1 & \text { при } & x>b, a<b .

\end{array}\right.

\]



В частнөм случае $a=-\infty(b=\infty)$ У. р. наз. у с еч ен ны м с п р а в а (с л е в a).



Наряду с (1) рассматриваются У. р. вида



\[

F_{a, b}(x)=\left\{\begin{array}{cll}

0 & \text { при } & x \leqslant a, \\

F(x)-F(a) & \text { при } & a<x<c, \\

F(x)+1-F(b) & \text { при } & c \leqslant x<b, \\

1 & \text { прш } & x \geqslant b,

\end{array}\right.

\]



Или



\[

F_{a, b}(x)=\left\{\begin{array}{ccl}

0 & \text { при } & x<a, \\

F(x) & \text { при } & a \leqslant x<b, \\

1 & \text { при } & x=b .

\end{array}\right.

\]



В (1) масса, сосредоточенная вне $[a, b]$, распределяется по всему отрезку $[a, b]$, в (2) - помещается в точку $c \in(a, b]$ (в том случае, когда $a<0<b$ в качестве $c$ чаще всего берется точка $c=0)$, а в (3) эта масса помеццается в крайние точки $a$ и $b$.



У. р. вида (1) могут интершретироваться следующим образом. Пусть $X$ - случайная величина с функцией распределения $F(x)$. Тогда У. р. совпадает с условным распределешием случаӥіпої величины при условии $a<$ $<X \leqslant b$.



С понятием У. р. тесно связано понятие усеченной случайной величины: если $X-$ случайная величина, то усеченной случайной величиной наз. величина



\[

X^{c}=\left\{\begin{array}{lll}

X, & \text { если } & |X| \leqslant c, \\

0, & \text { если } & |X|>c .

\end{array}\right.

\]



Распределение $X^{c}$ является У. р. тиша (3) по отношению к распределению $X$. Операция усечения - переход к У. р. или усеченным случайным величинам - является весьма распространенным технич. приемом. Она позволяет, «немного» изменяя исходное распределение, получить аналитич. своӥство - наличие всех моментов.



Лum.: [1] П p о $\mathrm{x}$ о $\mathrm{p}$ о в Ю. В., Р о 3 а н о в Ю. А., Теория вероятностей, 2 изд., М., 1973, [2] К р а м е р Г., Математические методы статистики, пер. с англ, 2 изд., М, 1975; жения, пер. с ангі., 2 изд., т. 2, М., 1967, [4] ЈІ о ә в М , Теория вероятностей, пер. с анг.і., М., 1962 Н. Г.' Ущаков.



УСЛОВНАЯ ВЕРОЯТНОСТЬ - 1) У. в. о т н о с ит е ль н о с о б ы т и я - характеристика связи двух событиї. Если $A$ и $B$ - события п $P(B)>0$, то У.в. $P(A \mid B)$ события $A$ относительно (или при условии) $B$ определяется равенством



\[

P(A \mid B)=\frac{P(A \cap B)}{P(B)} .

\]



У. в. $P(A \mid B)$ может рассматриваться как вероятность осуществления события $A$ при условии, что событие $B$ осупцестилось. Для независимых событиӥ $A$ и $B$ У. в. $P(A \mid B)$ совпадает с безусловной вероятностью $P(A)$.



О связи условных и безусловных вероятностей событий см. Бейеса формула и Полной вероятиости формула.



\begin{enumerate}

  \setcounter{enumi}{1}

  \item У. в. с обытия $A$ относ, т тельн о б-а лг е бр ы $\mathfrak{B}-$ случайная величина $P(A \mid \mathfrak{B})$, измеримая относительно $\mathfrak{B}$, для к-рой

\end{enumerate}



\[

\int_{B} P(A \mid \mathfrak{B}) P(d(1))=P(A \cap B)

\]



при любом $B \in \mathfrak{B}$. У. в. относительно б-алгебры определяется с точностью до эквивалентності.



Если б-алгебра $\mathfrak{B}$ порождена счетным числом непересекающихся событий $B_{1}, B_{2}, \ldots$, имеющих положительные вероятности и в сумме составляющих все пространство $\Omega$, то



$P(A \mid \mathfrak{B})=P\left(A \mid B_{k}\right) \quad$ при $\quad \omega \in B_{k}, \quad k-1,2, \ldots$



У. в. события $A$ относительно б-алгебры $\mathfrak{B}$ может быть определена как условиое математическое ожидание $\mathrm{E}\left(I_{A} \mid \mathfrak{B}\right)$ Індикатора $A$.



Пусть $(\Omega, \mathcal{A}, \quad$ ) - вероятностное пространство, $\mathfrak{B}$-под-б-алгебра $\mathcal{A}$. У. в. $P(A \mid \Re)$ наз. р е г у ли ян о й, если существует функция $p(\omega, A), \omega \in \Omega, A \in \mathcal{A}$, такая, что



a) при фиксированном $\omega$ фуннцця $p(\omega, A)$ является вероятностью на $\sigma$-алгебре $\mathscr{A}$,



б) $P(A \mid \mathfrak{B})=p(\omega, \boldsymbol{A})$ с вероятностью единица.



Для регулярных У. в. условные математит. ожидания выражаются интегралами с, У. в. в качестве мер.



У. в. относительно случаїноіі величины $X$ определяется как У. в. относительно б-алгебры, порожденноӥ $X$



Лит [1] К о л м огор о в А Н Основиые полятия теории вероятностей, 2 изд., М., 1974, 22$]$ П р о х ор о в Ю. В., Р о з а н о в 10. А., Теория вероятностей, 2 изд., М., 1973, [3] Л о ә в М., Терория вероятностей, пер. с аюг., М., 1962.



УСЈОВНАЯ ПЛОТНОСТЬ-плотность условного распределения. ІІусть $(\Omega, \mathcal{A}, \mathrm{P})$-вероятностное пространство, $\Re$ есть б-алгебра борелевских множкест на пря-



\begin{center}

\includegraphics[max width=\textwidth]{2023_07_09_3fe7e377419f8ee4ddc1g-1}

\end{center}



\[

Q(\omega, B)=\mathrm{P}\{X \in B \mid F\}, \quad \omega \in \Omega, \quad B \in \mathfrak{B},

\]



\begin{itemize}

  \item условное распределение $X$ относительно б-алгебры

\end{itemize}



\[

F_{X}(x \mid F)=Q(\omega,(-\infty, x))

\]



\begin{itemize}

  \item условная фуунция распределения $X$ относительно ЕслII

\end{itemize}



\[

F_{X}(x \mid F)=\int_{-\infty}^{x} f_{X}(t \mid F) d t

\]



то $f_{X}(x \mid F)$ наз. у сло вно й п лотн о с тью распределения $X$ относительно $\sigma$-алгебры





\end{document}